%% kapa: exam=2 round=29 year=2018 period=3 subject=compensation_law subject_ko=감정평가 및 보상법규 seq=3 number=3 points=20
%% source: Q-Net 감정평가사 기출문제 (cst00309) artlSeq=5203398 file=raw/1008/5203398/3교시 감정평가 및 보상법규.pdf
%% conversion: kordoc 4.24.51 -> tools/gen2.py (auto); review=proofread-gpt-terra; flags=none
\begin{longquestion}{【문제 3】}
서울의 A구청장은 이 사건 B토지의 비교표준지로 A구의 C토지(2017. 1. 1. 기준 공시지가는 1 m\textsuperscript{2}당 810만원임)를 선정하고 이 사건 B토지와 비교표준지 C의 토지가격비준표상 토지특성을 조사한 결과 총 가격배율이 1.00으로 조사됨에 따라 이 사건 각 토지의 가격을 1 m\textsuperscript{2}당 810만원으로 산정하였다. 감정평가사 D는 A구청장으로부터 이와 같이 산정된 가격의 검증을 의뢰받고 이 사건 각 토지가 비교표준지와 비교하여 환경조건, 획지조건 및 기타조건에서 열세에 있어 비교표준지의 공시지가를 약 83.9\%의 비율로 감액한 1 m\textsuperscript{2}당 680만원을 개별공시지가로 정함이 적정하다는 검증의견을 제시하였다. A구청장은 이 검증의견을 받아들여 2017. 5. 30.에 이 사건 각 토지의 개별공시지가를 1 m\textsuperscript{2}당 680만원으로 결정ㆍ공시하였다. B토지소유자는 1 m\textsuperscript{2}당 680만원으로 결정ㆍ공시된 B토지의 개별공시지가에 대하여 1 m\textsuperscript{2}당 810만원으로 증액되어야 한다는 취지로 이의신청을 제기하였다. B토지소유자의 이의신청에 따라 A구청장은 감정평가사 E에게 이 사건 토지의 가격에 대한 검증을 의뢰하였다. 검증을 담당한 감정평가사 E는 토지특성 적용 및 비교표준지 선정에는 오류가 없으나 인근 지가와의 균형을 고려하여 개별공시지가를 1 m\textsuperscript{2}당 700만원으로 증액함이 상당하다는 의견을 제시하였다. (이 사건 토지가 비교표준지와 비교하여 환경조건 및 획지조건에서 열세에 있다고 보아 비교표준지의 공시지가에 대하여 약 86.5\%의 비율로 감액) \par
이에 A구청장은 A구 부동산가격공시위원회의 심의를 거쳐 이 검증의견을 받아들여 B토지에 대하여 1 m\textsuperscript{2}당 700만원으로 개별공시지가결정을 하였다. 이에 대하여 B토지소유자는 토지가격비준표와 달리 결정된 개별공시지가결정은 위법하다고 주장한다. 이 주장은 타당한가? (20점) \par
\begin{framedtext}
참조 조문 \par
<부동산 가격공시에 관한 법률> 
제10조(개별공시지가의 결정ㆍ공시 등) ① 시장ㆍ군수 또는 구청장은 국세ㆍ지방세 등 각종 세금의 부과, 그 밖의 다른 법령에서 정하는 목적을 위한 지가산정에 사용되도록 하기 위하여 
제25조에 따른 시ㆍ군ㆍ구부동산가격공시위원회의 심의를 거쳐 매년 공시지가의 공시기준일 현재 관할 구역 안의 개별토지의 단위면적당 가격(이하 "개별공시지가"라 한다)을 결정ㆍ공시하고, 이를 관계 행정기관 등에 제공하여야 한다. 
②\textasciitilde{}③ 생략 
④ 시장ㆍ군수 또는 구청장이 개별공시지가를 결정ㆍ공시하는 경우에는 해당 토지와 유사한 이용가 치를 지닌다고 인정되는 하나 또는 둘 이상의 표준지의 공시지가를 기준으로 토지가격비준표를 사용하여 지가를 산정하되, 해당 토지의 가격과 표준지공시지가가 균형을 유지하도록 하여야 한다. 
⑤ <이하 생략> <부동산 가격공시에 관한 법률 시행령> 
제18조(개별공시지가의 검증) ① 생략 
② 법 
제10조제5항 본문에 따라 검증을 의뢰받은 감정평가업자는 다음 각 호의 사항을 검토ㆍ확인하고 의견을 제시하여야 한다. 
1. 비교표준지 선정의 적정성에 관한 사항 
2. 개별토지 가격 산정의 적정성에 관한 사항 
3. 산정한 개별토지가격과 표준지공시지가의 균형 유지에 관한 사항 
4. 산정한 개별토지가격과 인근토지의 지가 및 전년도 지가와의 균형 유지에 관한 사항 
5. 그 밖에 시장ㆍ군수 또는 구청장이 검토를 의뢰한 사항
\end{framedtext}

\end{longquestion}
